% !TeX root = ../../Book1.tex
% 纲要标题：### F.1 单项式序与多元除法
% 纲要位置：工作记录/计划纲要.md，第 955 行。

% 纲要要点（供目录核对）：
% #### F.1.1 单项式的整除与次序
% * 多重指数与单项式整除，单项式的最大公因式和最小公倍式。
% * 全局单项式序的全序、良序和乘法相容条件。
% * 以字典序、分次字典序为主要例子，比较它们选择首项的方式；分次逆字典序只作简短补充。
% #### F.1.2 首项与多元除法
% * 首单项式、首项系数和首项的区别；固定记号与零多项式的处理约定。
% * 写出除以一组多项式的逐步算法，保存等式 \(f=\sum_iq_ig_i+r\)。
% * 用首单项式的严格下降证明终止，说明余式中的单项式均不能被任何除式的首单项式整除。
% #### F.1.3 余式为何依赖除式顺序
% * 在 \(\mathbb Q[x,y]\) 中固定字典序 \(x>y\)，用 \(xy-1,y^2-1\) 除 \(xy^2\)，分别得到余式 \(y\) 与 \(x\)。
% * 解释零余式总能提供理想成员证据，而对任意生成组得到非零余式，不能据此断定不属于理想。
% * 引出问题：怎样选取生成组，才能让余式唯一并准确判断理想成员？

\section{单项式序与多元除法}
\label{b1:sec:F01}

一元多项式除法每次消去次数最高的一项，次数便严格下降。到了两个变量，\(x^2\) 与 \(y^3\) 哪一项应当先消去，却不能仅由“多项式”这个概念决定。必须先为单项式选定次序，再问当前首项能否被某个除式的首项整除。这样的除法可以完成，但它最初并不具有一元除法的全部优点：同一个被除式，换一个除式顺序，就可能得到不同的余式。

\ProseParagraphBreak
本附录始终令 \(k\) 为域、\(n\geq1\) 为整数，并置 \(R=k[x_1,\ldots,x_n]\)。一般结论不限制 \(k\) 的特征；实际运行算法时，还要求能够在 \(k\) 中进行精确四则运算并判定一个系数是否为零。有理数域和给定的有限域满足这些要求。

\subsection{单项式的整除与次序}
\label{b1:subsec:F01:01}

沿用\cref{b1:def:multivariate-degree}前的记号，若 \(\alpha=(\alpha_1,\ldots,\alpha_n)\in\mathbb N^n\)，则
\[
x^\alpha=x_1^{\alpha_1}\cdots x_n^{\alpha_n},\qquad
|\alpha|=\alpha_1+\cdots+\alpha_n.
\]
这里单项式的系数取为 \(1\)；\(c x^\alpha\)，其中 \(c\in k\) 非零，称为多项式的一项。单项式的整除可以逐个变量检查：
\[
x^\alpha\mid x^\beta
\quad\Longleftrightarrow\quad
\alpha_i\leq\beta_i\quad(1\leq i\leq n).
\]
在满足这些不等式时，商是 \(x^{\beta-\alpha}\)。因此最大公因式和最小公倍式分别由各分量的最小值和最大值给出：
\[
\begin{aligned}
\gcd(x^\alpha,x^\beta)
 &=\prod_{i=1}^n x_i^{\min(\alpha_i,\beta_i)},\\
\operatorname{lcm}(x^\alpha,x^\beta)
 &=\prod_{i=1}^n x_i^{\max(\alpha_i,\beta_i)}.
\end{aligned}
\]
例如 \(x^2y\) 与 \(xy^3\) 的最大公因式是 \(xy\)，最小公倍式是 \(x^2y^3\)。整除只是偏序：\(x\) 与 \(y\) 互不整除；消项算法需要另外选一个能比较任意两个单项式的次序。

\begin{definition}\label{b1:def:F-monomial-order}
设 \(k\) 为域，\(n\geq1\) 为整数，\(R=k[x_1,\ldots,x_n]\)。若 \(R\) 的单项式集合上的全序 \(\preceq\) 满足以下条件：
\begin{enumerate}
\item 任一非空单项式集合都有最小元，即这个全序是良序；
\item 对任意单项式 \(u,v,w\)，若 \(u\prec v\)，则 \(uw\prec vw\)。
\end{enumerate}
则称 \(\preceq\) 为 \(R\) 上的\term[danxiangshi-xu]{单项式序}[monomial order]。本附录只采用这样的全局单项式序；符号 \(u\prec v\) 表示 \(u\preceq v\) 且 \(u\ne v\)。
\end{definition}

这些条件蕴含 \(1\) 是最小单项式。否则若 \(u\prec1\)，乘法相容性就给出无穷严格下降链 \(1\succ u\succ u^2\succ\cdots\)，其各项组成的集合没有最小元，与良序矛盾。因而只要 \(u\mid v\)，就有 \(u\preceq v\)；若 \(u\ne v\)，这个不等式严格成立。反过来，两个单项式在所选次序下可比，并不意味着它们存在整除关系。

\ProseParagraphBreak
最直接的例子是\term[zidian-xu]{字典序}[lexicographic order]。规定变量优先级为 \(x_1\succ\cdots\succ x_n\)，比较 \(x^\alpha\) 和 \(x^\beta\) 时，从第一个坐标开始找出 \(\alpha_j\ne\beta_j\) 的位置，并规定
\[
x^\alpha\prec_{\mathrm{lex}}x^\beta
\quad\Longleftrightarrow\quad \alpha_j<\beta_j.
\]
这与按首个不同字母排列单词相似，但这里比较的是各变量的指数。它确实是良序：在任一非空指数集合中，先把第一坐标限制为其中的最小值，再在剩下的集合中把第二坐标限制为最小值，经过有限个坐标后得到字典序最小的指数。给两组指数同时加上 \(\gamma\)，首个不同位置及其大小关系不变，所以乘法相容性也成立。字典序不按全次数排列；在 \(k[x,y]\) 中取 \(x\succ y\) 时，对任意 \(m\geq1\) 都有 \(x\succ y^m\)。

\ProseParagraphBreak
若希望先比较全次数，可以采用\term[fenci-zidian-xu]{分次字典序}[graded lexicographic order]：先比较 \(|\alpha|\) 和 \(|\beta|\)，全次数相同时再按字典序比较。任何非空单项式集合先有最小全次数，而固定全次数的单项式只有有限多个，因此这是良序。乘以同一个单项式时，两边全次数同时增加同一个数；若原来全次数相等，字典序的比较又保持不变。所以它也是单项式序。例如 \(x+y^2\) 在字典序 \(x\succ y\) 下以 \(x\) 为最大单项式，在分次字典序下却以 \(y^2\) 为最大单项式。

\ProseParagraphBreak
计算中还常用\term[fenci-ni-zidian-xu]{分次逆字典序}[graded reverse lexicographic order]：先比较全次数；全次数相同时，从最后一个坐标向前寻找首个不同位置 \(j\)，并规定指数较小者的单项式较大，即 \(\alpha_j<\beta_j\) 时 \(x^\alpha\succ x^\beta\)。在每一固定全次数上，这就是对反向排列并改变符号的指数作字典序比较，所以是全序。对非空单项式集合，先选其中最小的全次数，再在该次数的有限集合中取最小元，便得到良序性；加上同一个指数向量仍保持比较，故乘法相容。例如在 \(k[x,y,z]\) 中，\(x^2z\) 与 \(xy^2\) 的全次数都为三：字典序和分次字典序把前者排得更大，分次逆字典序则把后者排得更大。后文证明只使用\cref{b1:def:F-monomial-order}的两项条件；每个算例则会明确所取次序。

\subsection{首项与多元除法}
\label{b1:subsec:F01:02}

固定一个单项式序。非零多项式只含有限多个单项式，因而其中有唯一的最大者。把它与所带系数分开记，可以避免在整除和消项时混淆两类数据。

\begin{definition}\label{b1:def:F-leading-data}
设 \(k\) 为域，\(n\geq1\) 为整数，\(R=k[x_1,\ldots,x_n]\) 配备一个单项式序。对非零多项式 \(f=\sum_\alpha c_\alpha x^\alpha\in R\)，设 \(x^\beta\) 是满足 \(c_\beta\ne0\) 的单项式中最大者。令
\[
\operatorname{LM}(f)=x^\beta,\qquad
\operatorname{LC}(f)=c_\beta,\qquad
\operatorname{LT}(f)=c_\beta x^\beta
\]
分别称 \(\operatorname{LM}(f)\)、\(\operatorname{LC}(f)\) 与 \(\operatorname{LT}(f)\) 为 \(f\) 的\term[shou-danxiangshi]{首单项式}[leading monomial]、首项系数与\term[shouxiang]{首项}[leading term]。这些记号总是相对于当前固定的单项式序。零多项式不定义首单项式和首项系数；含有这些记号的比较均只针对非零多项式。
\end{definition}

\termidx[shouxiang-xishu]{首项系数}
\symidx{\(\operatorname{LM}(f)\)：多项式的首单项式}
\symidx{\(\operatorname{LC}(f)\)：多项式的首项系数}
\symidx{\(\operatorname{LT}(f)\)：多项式的首项}

例如在 \(\mathbb Q[x,y]\) 中，取字典序 \(x\succ y\)，则 \(f=3x^2y-5xy^4+2\) 满足 \(\operatorname{LM}(f)=x^2y\)、\(\operatorname{LC}(f)=3\)、\(\operatorname{LT}(f)=3x^2y\)；若换为分次字典序，三者依次变成 \(xy^4\)、\(-5\)、\(-5xy^4\)。因此“首”表示所选单项式序中的最大位置，不总是全次数最高的位置。

\begin{proposition}\label{b1:prop:F-leading-product}
设 \(k\) 为域，\(n\geq1\) 为整数，\(R=k[x_1,\ldots,x_n]\) 配备一个单项式序。若 \(f,g\in R\) 均非零，则
\[
\operatorname{LM}(fg)=\operatorname{LM}(f)\operatorname{LM}(g),\qquad
\operatorname{LT}(fg)=\operatorname{LT}(f)\operatorname{LT}(g).
\]
若 \(f+g\ne0\)，则
\[
\operatorname{LM}(f+g)\preceq
\max\Bigl\{\operatorname{LM}(f),\operatorname{LM}(g)\Bigr\}.
\]
当 \(f\)、\(g\) 的首单项式不同，和式中的不等式取等号。
\end{proposition}
\begin{proof}
记 \(u=\operatorname{LM}(f)\)、\(v=\operatorname{LM}(g)\)。若 \(u'\)、\(v'\) 分别来自 \(f\)、\(g\)，则 \(u'\preceq u\)、\(v'\preceq v\)，所以 \(u'v'\preceq uv\)。除 \((u',v')=(u,v)\) 外，至少一个比较严格；乘法相容性使乘积比较也严格。因此 \(uv\) 只来自两个首项的乘积，其系数是两个非零首项系数的积，在域中不为零。这就证明前两式。相加不会引入原来两式都没有的单项式；若两个首单项式不同，较大者也没有另一项可以与它抵消，故得到最后的断言。
\end{proof}

现在给定有序的一组非零多项式 \(g_1,\ldots,g_s\)。若当前多项式 \(p\) 的首单项式被 \(\operatorname{LM}(g_i)\) 整除，就可以减去
\[
\frac{\operatorname{LT}(p)}{\operatorname{LT}(g_i)}g_i
=\frac{\operatorname{LC}(p)}{\operatorname{LC}(g_i)}
 \frac{\operatorname{LM}(p)}{\operatorname{LM}(g_i)}g_i
\]
来消去首项。右边两个分式分别是域中的系数商和整除成立时的单项式商；并没有在 \(R\) 中对任意多项式作除法。

\ProseParagraphBreak
把这一步反复应用，就得到下面的\term[duoyuan-chufa]{多元除法}[multivariate division]。若当前首项不能消去，算法并不立刻停下，而是把这一项移入余式，再处理剩余部分。

\begin{theorem}[多元除法]\label{b1:thm:F-division}
设 \(k\) 为域，\(n\geq1\) 为整数，\(R=k[x_1,\ldots,x_n]\) 配备一个单项式序。对 \(f\in R\) 和有限个非零多项式 \(g_1,\ldots,g_s\in R\)，可以经过有限次消项求出 \(q_1,\ldots,q_s,r\in R\)，使
\begin{equation}\label{b1:eq:F-division-representation}
f=q_1g_1+\cdots+q_sg_s+r,
\end{equation}
且 \(r\) 中的每个单项式都不能被任何 \(\operatorname{LM}(g_i)\) 整除。若 \(f\ne0\)，输出还满足：每个非零的 \(q_i g_i\) 以及非零的 \(r\)，其首单项式都不大于 \(\operatorname{LM}(f)\)。若 \(f=0\)，可取全部输出为零；若除式组为空，则取 \(r=f\)。
\end{theorem}
\begin{proof}
从 \(p=f\)、\(q_1=\cdots=q_s=r=0\) 开始，始终保持等式
\begin{equation}\label{b1:eq:F-division-invariant}
f=q_1g_1+\cdots+q_sg_s+r+p.
\end{equation}
只要 \(p\ne0\)，就按下列步骤处理。
\begin{enumerate}
\item 若有 \(i\) 满足 \(\operatorname{LM}(g_i)\mid\operatorname{LM}(p)\)，取符合条件的最小指标 \(i\)，令 \(t=\operatorname{LT}(p)/\operatorname{LT}(g_i)\)，把 \(q_i\) 改为 \(q_i+t\)，把 \(p\) 改为 \(p-tg_i\)。首项被消去；由\cref{b1:prop:F-leading-product}，其余留下的单项式均严格小于原来的 \(\operatorname{LM}(p)\)。
\item 若没有这样的 \(i\)，把 \(r\) 改为 \(r+\operatorname{LT}(p)\)，把 \(p\) 改为 \(p-\operatorname{LT}(p)\)。这一项无法再用任何除式的首项消去，所以移入余式。新的非零 \(p\) 的首单项式仍严格减小。
\end{enumerate}
两种操作都保持\eqref{b1:eq:F-division-invariant}。若过程永不结束，各个非零 \(p\) 的首单项式就形成无穷严格下降链，与单项式序的良序性矛盾。因此有限步后 \(p=0\)，得到\eqref{b1:eq:F-division-representation}。每个移入 \(r\) 的单项式都不被任何 \(\operatorname{LM}(g_i)\) 整除；相加至多消去已有的项，不会产生新的单项式，所以最终 \(r\) 仍满足这一条件。

最后检查首项界。每次加入某个 \(q_i\) 的项 \(t\) 都满足 \(\operatorname{LM}(tg_i)=\operatorname{LM}(p)\preceq\operatorname{LM}(f)\)，而移入 \(r\) 的项也满足同一上界。把这些项分别相加不会产生更大的单项式；由\cref{b1:prop:F-leading-product}，所有非零的 \(q_i g_i\) 和非零的 \(r\) 都满足所述上界。
\end{proof}

算法可以选择任意一个能够消项的除式，终止性和所述首项界都不变。规定“最小指标”只是把选择固定下来，使同一有序输入产生确定的输出。域的假设用在首项系数求逆处；良序性负责终止，而乘法相容性保证一次消项不会反而产生更大的单项式。这三件事在论证中各有不同作用。

\subsection{余式为何依赖除式顺序}
\label{b1:subsec:F01:03}

在 \(\mathbb Q[x,y]\) 中固定字典序 \(x\succ y\)，取
\[
g_1=xy-1,\qquad g_2=y^2-1,\qquad f=xy^2.
\]
\(f\) 的首单项式同时被 \(\operatorname{LM}(g_1)=xy\) 和 \(\operatorname{LM}(g_2)=y^2\) 整除。如果先用 \(g_1\)，就有
\[
xy^2-y(xy-1)=y,
\qquad f=yg_1+y.
\]
\(y\) 不能再被 \(xy\) 或 \(y^2\) 整除，所以余式为 \(y\)。如果先用 \(g_2\)，则
\[
xy^2-x(y^2-1)=x,
\qquad f=xg_2+x,
\]
这次余式为 \(x\)。两次计算均符合\cref{b1:thm:F-division}，但结果不同。差异不是计算失误，也不是系数近似所致；它来自允许消项的先后选择。

\ProseParagraphBreak
将两条表示相减，得到
\begin{equation}\label{b1:eq:F-missing-generator}
x-y=yg_1-xg_2\in(g_1,g_2).
\end{equation}
然而 \(x-y\) 的两项都不能被 \(xy\) 或 \(y^2\) 整除。因此用这两个除式去除 \(x-y\)，算法不消去任何一项，得到的非零余式就是 \(x-y\) 自身。对于任意给定的生成组，“余式非零”因而不能直接推出被除式不在所生成的理想中。

\ProseParagraphBreak
零余式的意义则没有这种不确定性。如果一次除法得到 \(r=0\)，\eqref{b1:eq:F-division-representation}就是 \(f\in(g_1,\ldots,g_s)\) 的具体证据，并记录了所需的系数 \(q_i\)。真正缺少的是反向判据：当 \(f\) 属于理想时，能否保证除法必得零余式？例子中的 \(x-y\) 提示了改进办法：原生成组没有直接显露理想中每一种可能的首项；把必要的关系补进去，才可能使消项过程准确反映整个理想。下一节把这一要求写成精确定义。
